% English translation of capitulos/07-apendice-busca.tex, made by Claude Code (Claude Fable 5.2) on 24 November 2026.
% Pending review by the author. When the Portuguese file changes, update this one (scripts/verificar_traducao.py).
\appendices
\section{Search string and filters}\label{ap:busca}
The general search string, applied to title and abstract, is:

{\footnotesize\raggedright
(``prompt engineering'' OR prompting OR prompt OR prompts OR ``in-context learning'' OR ``few-shot''
OR ``zero-shot'' OR ``chain-of-thought'' OR ``chain of thought'') AND (``code generation'' OR
``program synthesis'' OR ``code translation'' OR ``program translation'' OR ``code migration'' OR
transpil* OR refactoring OR refactor OR debugging OR debug OR ``program repair'' OR ``bug fixing'' OR
``bug fix'' OR ``bug repair'') AND (``large language model'' OR ``large language models'' OR LLM OR
LLMs OR ``language model'' OR ``language models'' OR ``generative AI'' OR ChatGPT OR ``GPT-4'' OR
Codex)\par}

\medskip
The filters applied were: on arXiv, submission date from 1 January 2021 to 7 September 2026; on IEEE Xplore,
years from 2021 to 2026; on the ACM Digital Library, electronic publication from January 2021 to September 2026, including the
whole month, in the ACM Full-Text collection. In the rapid review adopted, the block of tasks is applied only to the
title, by local simulation over the records, without a new search in the databases. The exact query for each database is in \texttt{dados/buscas.csv}, in the project repository.
